% -*- coding: utf-8 -*-
% !TEX program = xelatex
\documentclass[plain,noanswer,medmath]{../../template/jnuexam}
\usepackage{../../template/kysx}

\begin{document}


\examtitle{2006}{2}


\loadproblems{1}{2006P1}%载入试卷一


\makepart{填空题}{1～6小题，每小题4分，共24分}


\begin{problem}
曲线$y = \frac{x + 4\sin x}{5x - 2\cos x}$的水平渐近线方程为 \fillin{}．
\end{problem}


\begin{problem}
设函数$f(x) = \begin{cases}
  \frac{1}{x^3}\int_0^x \sin t^2 \dt, & x \ne 0 \\
                                   a, & x = 0
\end{cases}$在$x = 0$处连续，则$a =$ \fillin{}．
\end{problem}


\begin{problem}
广义积分$\int_0^{+\infty}\!\!\frac{x\dx}{(1 + x^2)^2} =$ \fillin{}．
\end{problem}


\useproblem{1}{1}{2}%【同试卷一第一(2)题】


\begin{problem}
设函数$y = y(x)$由方程$y = 1 - x\e^y$确定，则$\frac{\dy}{\dx}\Big|_{x=0}=$ \fillin{}．
\end{problem}


\useproblem{1}{1}{5}%【同试卷一第一(5)题】


\makepart{选择题}{7～14小题，每小题4分，共32分}


\useproblem{1}{2}{7}%【同试卷一第二(7)题】


\begin{problem}
设$f(x)$是奇函数，除$x = 0$外处处连续，$x = 0$是其第一类间断点，则$\int_0^x f(t)\dt$是\pickout{}
\begin{abcd}
\item 连续的奇函数．
\item 连续的偶函数．
\item 在$x = 0$间断的奇函数．
\item 在$x = 0$间断的偶函数．
\end{abcd}
\end{problem}


\begin{problem}
设函数$g(x)$可微，$h(x) = \e^{1 + g(x)}$，$h'(1) = 1$，$g'(1) = 2$，则$g(1)$等于\pickout{}
\begin{abcd}
\item $\ln 3 - 1$．
\item $- \ln 3 - 1$．
\item $- \ln 2 - 1$．
\item $\ln 2 - 1$．
\end{abcd}
\end{problem}


\begin{problem}
函数$y = c_1\e^x + c_2\e^{- 2x} + x\e^x$满足的一个微分方程是\pickout{}
\begin{abcd}
\item $y'' - y' - 2y = 3x\e^x$．
\item $y'' - y' - 2y = 3\e^x$．
\item $y'' + y' - 2y = 3x\e^x$．
\item $y'' + y' - 2y = 3\e^x$．
\end{abcd}
\end{problem}


\useproblem{1}{2}{8}%【同试卷一第二(8)题】


\useproblem{1}{2}{10}%【同试卷一第二(10)题】


\useproblem{1}{2}{11}%【同试卷一第二(11)题】


\useproblem{1}{2}{12}%【同试卷一第二(12)题】


\makepart{解答题}{15～23小题，共94分}


\begin{problem}[points=10]
试确定常数$A,B,C$的值，使得$e^x(1 + Bx + Cx^2) = 1 + Ax + o(x^3)$，
其中$o(x^3)$是当$x \to 0$时比$x^3$高阶的无穷小．
\end{problem}


\begin{problem}[points=10]
求$\int \frac{\arcsin \e^x}{\e^x}\dx$．
\end{problem}


\useproblem[points=10]{1}{3}{15}%【同试卷一第三(15)题】


\useproblem[points=12]{1}{3}{16}%【同试卷一第三(16)题】


\begin{problem}[points=10]
证明：当$0 < a < b < \pi$时，$b\sin b + 2\cos b + \pi b > a\sin a + 2\cos a + \pi a$．
\end{problem}


\useproblem[points=12]{1}{3}{18}%【同试卷一第三(18)题】


\begin{problem}[points=12]
已知曲线$L$的方程$\begin{cases}
  x = t^2 + 1, \\
  y = 4t - t^2,
\end{cases}$（$t \ge 0$）．
\begin{enumerate}
\item 讨论$L$的凹凸性；
\item 过点$(-1,0)$引$L$的切线，求切点$(x_0,y_0)$，并写出切线的方程；
\item 求此切线与$L$（对应$x \le x_0$的部分）及$x$轴所围成的平面图形的面积．
\end{enumerate}
\end{problem}


\useproblem[points=9]{1}{3}{20}%【同试卷一第三(20)题】


\useproblem[points=9]{1}{3}{21}%【同试卷一第三(21)题】


\end{document}
